% ============================================================
\section{Iteración 1 --- DRV-01: servidor autoritativo}\label{sec:it01}
% ============================================================

% ------------------------------------------------------------
\subsection{Driver y alcance}

\begin{tabla}{|L{3.0cm}|Y|}
\fichacab
\parte{Driver}{DRV-01 --- Servidor autoritativo.}
\parte{Enunciado}{El servidor arbitra cada partida en línea y es la única autoridad sobre su estado: el cliente propone y presenta; el servidor decide y conserva el estado vigente, con independencia de las conexiones y de la base de datos.}
\parte{Por qué es la primera}{Obtiene la calificación máxima en los cuatro criterios de la sección~\ref{sec:priorizacion}: agrupa cuatro ASR prioritarios, entre ellos ASR-08 y ASR-02, que abren el orden de la sección~\ref{sec:asr}; cambia cliente, servidor, protocolo y capa de datos a la vez; es la decisión que CRN-10 llama la más cara de cambiar; y habilita a DRV-02, DRV-03, DRV-04 y DRV-06.}
\parte{Qué decide esta iteración}{Quién produce y conserva el estado vigente de una partida en curso, y dónde vive mientras la partida dura. Nada más: es el único objetivo del paso~2 y se comprueba con una sola medida, la de ESC-05.}
\parte{Decisiones ya tomadas}{El equipo ya fijó el modelo cliente-servidor basado en servicios, que no se confía en el cliente y que toda la lógica del juego se ejecuta en el servidor, y que una partida en línea exige conexión con la base de datos. Esta iteración no las vuelve a discutir: las toma como entradas y desarrolla sus consecuencias (DEC-01 a DEC-03, registradas en la sección~\ref{sec:asr}).}
\parte{Qué no decide}{Las otras dos medidas de DRV-01 tienen iteración propia porque dependen de otros elementos: la supervivencia a una desconexión (ESC-08) se decide en una iteración posterior, sobre la sesión de partida, y el comportamiento ante una caída de la base de datos (ESC-09) en la iteración~5, sobre el servicio de acceso a datos. Quedan fuera también el protocolo (iteración~2), el motor de reglas (iteración~3) y el reparto del segundo (iteración~4). No resuelve Q-01, Q-02 ni Q-09.}
\end{tabla}

% ------------------------------------------------------------
\subsection{Paso 1. Revisión de entradas}

Se revisan las cinco entradas del método y se listan las que tienen algo que ver con el driver. Las demás siguen siendo válidas, pero no dicen nada sobre dónde vive el estado de una partida.

\subsubsection*{ASR que agrupa el driver}

Del conjunto de la sección~\ref{sec:asr}, estos son los ASR que DRV-01 agrupa y el papel que cumple cada uno en esta iteración.

\begin{tabla}{|L{1.4cm}|Y|L{1.7cm}|L{3.6cm}|}
\hline
\filacab \cab{ASR} & \cab{Qué debe cumplir la arquitectura} & \cab{Prioridad} & \cab{Papel en esta iteración} \\ \hline
\endhead
ASR-02 & Un cliente alterado no puede cambiar el estado de una partida ni obtener datos de otro jugador. & Prioritario & Objetivo. Su medida es ESC-05. \\ \hline
ASR-08 & Dos o más jugadores juegan partidas completas en línea, creadas por emparejamiento o por sala, con turnos y relojes arbitrados por un servidor. & Prioritario & Base del incremento: la partida que CU-20 y CU-21 recorren de principio a fin. \\ \hline
ASR-04 & Una desconexión no ordenada no deja la partida indefinida: se detecta en menos de 15 segundos y el jugador retoma en menos de 5 si vuelve antes de 2 minutos, incluso desde otro dispositivo. & Prioritario & Pendiente: necesita el protocolo de la iteración~2. Aquí solo se evita atar el estado al socket. \\ \hline
ASR-05 & Una caída de SQL Server pierde como máximo 5 segundos de jugadas y ningún dato confirmado, y deja constancia de lo perdido. & Prioritario & Se atiende en la iteración~5. Aquí solo se decide que la base de datos no es la copia de trabajo. \\ \hline
ASR-09 & Una partida en curso se puede observar, incluso sin cuenta, recibiendo el estado con retraso y sin frenar a los jugadores. & Media & Condición en la iteración~4. \\ \hline
\end{tabla}

\subsubsection*{Objetivos de diseño}

\begin{tabla}{|L{1.5cm}|L{5.0cm}|Y|}
\hline
\filacab \cab{ID} & \cab{Objetivo} & \cab{Origen y qué exige aquí} \\ \hline
\endhead
OBJ-01 & Dos jugadores pueden jugar una partida completa en línea, de principio a fin. &
REQ-01. Es el objetivo que vuelve obligatorio tener un servidor: sin él no hay partida que arbitrar. \\ \hline
OBJ-02 & Nadie puede alterar una partida desde un cliente modificado. &
CRN-23 y el interés de STK-15. Da por hecho que el cliente es hostil, y eso decide quién tiene la última palabra sobre el estado. \\ \hline
OBJ-03 & La partida sobrevive a las fallas normales de una red doméstica y a una caída de la base de datos. &
CRN-06 y CRN-19, con STK-01, STK-13 y STK-14. Exige que el estado siga existiendo cuando se cae aquello a lo que estaba asociado. \\ \hline
OBJ-04 & Cada semana hay un incremento funcional y demostrable dentro de las 14 semanas. &
CON-13 y CON-14. Limita lo que esta iteración puede permitirse: la estructura tiene que poder construirse por partes y probarse desde el primer incremento. \\ \hline
\end{tabla}

\subsubsection*{Requisitos funcionales primarios}

Del documento de casos de uso se toman como primarios los que cambian el estado de una partida en curso o dependen de él. Son los que obligan al servidor a arbitrar y los que vuelven visible cualquier discrepancia entre lo que ve el jugador y lo que el servidor acepta.

\begin{tabla}{|L{1.3cm}|L{4.2cm}|Y|}
\hline
\filacab \cab{CU} & \cab{Caso de uso} & \cab{Qué exige del estado de la partida} \\ \hline
\endhead
CU-20 & Mover el peón &
Una de las dos acciones del turno. El servidor comprueba sesión, participación, turno, número de jugada y reloj antes de aceptarla, y notifica el estado actualizado a todos. \\ \hline
CU-21 & Colocar un muro &
La otra acción del turno y la más cara: exige comprobar que el muro no se solape ni se cruce (RN-04) y que todos los jugadores activos conserven camino a su meta (RN-05), sobre el tablero vigente. \\ \hline
CU-26 & Reconectar a una partida en curso &
El estado tiene que seguir vivo mientras el jugador no está, con su reloj corriendo, y devolverse completo al reconectar, incluso desde otro dispositivo (FA-05, MATCH\_TAKEN\_OVER). \\ \hline
CU-24 & Abandonar la partida &
Distingue entre irse y desconectarse: a los 60 segundos de desconexión el rival puede reclamar la victoria (RN-03 del CU-26), lo que exige que el tiempo lo mida el servidor y no el cliente. \\ \hline
CU-25 & Finalizar la partida y aplicar la progresión &
Es el único momento en que el resultado se vuelve definitivo: forma de término, resultado por participación, elo, progresión y monedas. Marca la frontera entre lo que se puede perder y lo que no. \\ \hline
CU-22, CU-23 & Ofrecer y aceptar tablas; rendirse &
Terminan la partida por acuerdo o por decisión de un jugador y dependen del mismo árbitro: una oferta pendiente caduca cuando llega una jugada (RN-05 del CU-22). \\ \hline
CU-17, CU-18, CU-19 & Buscar partida clasificatoria; crear sala privada; unirse con código &
Crean la partida y fijan sus parámetros. No se descomponen en esta iteración, pero determinan el momento en que el estado nace. \\ \hline
\end{tabla}

Quedan fuera de los primarios y se revisan en iteraciones posteriores: CU-27 (jugar contra la IA), que depende de Q-01; CU-28 y CU-29 (chat y emotes), que corresponden a DRV-06; CU-34 (ver una partida como espectador), que consume el estado con retraso pero no lo cambia; y CU-36 y CU-37 (historial y repetición), que dependen de lo que DRV-04 decida persistir.

\subsubsection*{Escenarios de atributos de calidad}

\begin{tabla}{|L{1.7cm}|L{3.3cm}|Y|}
\hline
\filacab \cab{Escenario} & \cab{Atributo} & \cab{Medida que esta iteración tiene que volver posible} \\ \hline
\endhead
ESC-05 \newline (ASR-02) & Seguridad: integridad &
Con 200 mensajes inválidos de un cliente modificado: el 100~\% se rechaza, el estado es idéntico antes y después, 0 respuestas con datos del rival, el servidor no se detiene y el 100~\% de los intentos queda registrado con cuenta y hora. \\ \hline
ESC-08 \newline (ASR-04) & Fiabilidad: recuperabilidad &
Detección de la desconexión en menos de 15 segundos; si el jugador vuelve antes de 2 minutos, retoma en menos de 5 segundos con el mismo tablero, el mismo turno y 0 jugadas perdidas; en 50 cortes, 0 partidas indefinidas. \\ \hline
ESC-09 \newline (ASR-05) & Fiabilidad y responsabilidad &
Tras una caída abrupta de SQL Server con 20 partidas en curso: 0 cuentas y progresos confirmados perdidos, como máximo las jugadas de los últimos 5 segundos, 0 registros a medias y ningún jugador esperando más de 10 segundos sin aviso. \\ \hline
ESC-01 \newline (ASR-01) & Desempeño &
No es el objetivo de esta iteración, pero sí su límite: ninguna decisión puede impedir que el 100~\% de 1\,000 jugadas responda en menos de 1 segundo, con mediana bajo 300 ms. \\ \hline
ESC-07 & Responsabilidad &
Exige que cada acción quede ligada a la cuenta que la hizo. Se apoya en el mismo registro que pide ESC-05. \\ \hline
\end{tabla}

\subsubsection*{Restricciones}

\begin{tabla}{|L{1.9cm}|L{4.3cm}|Y|}
\hline
\filacab \cab{ID} & \cab{Restricción} & \cab{Cómo limita esta iteración} \\ \hline
\endhead
REQ-01 & Modo multijugador en línea obligatorio & La partida no puede resolverse en un solo equipo: hay dos extremos y una red entre ellos. \\ \hline
CON-02 & La comunicación es por TCP & Hay entrega ordenada y confiable, pero también conexiones medio abiertas: una caída de Wi-Fi no cierra el socket. \\ \hline
CON-03 & Está prohibido usar WebSockets & Lo que un protocolo hereda normalmente lo construye el equipo. Ni la detección de actividad ni la reasociación de sesión vienen dadas. \\ \hline
CON-17 & La pila TCP/IP varía entre entornos & Impide confiar en el mecanismo del sistema operativo para detectar la caída: tarda horas y no es uniforme. \\ \hline
CON-04, DEP-03 & Almacenamiento en SQL Server, alojado en infraestructura externa & La base de datos puede caerse y su latencia no está bajo el control del equipo. \\ \hline
DEP-04 & Infraestructura de servidores y red provista externamente & El servidor puede reiniciarse sin aviso; la recuperación sí es responsabilidad del equipo. \\ \hline
CON-13, CON-14 & Plazo de 14 semanas e incrementos auditables & Descartan soluciones que no se puedan construir y probar por partes dentro del plazo. \\ \hline
CON-01, CON-16 & C\# sobre .NET & Acotan herramientas, no estructura. Abaratan que cliente y servidor compartan un mismo ensamblado. \\ \hline
\end{tabla}

\subsubsection*{Decisiones de proyecto ya tomadas}

Además de las restricciones que vienen del cliente, el equipo ya tomó tres decisiones que condicionan esta iteración. Entran como entradas y no como alternativas: lo que sigue desarrolla sus consecuencias.

\begin{tabla}{|L{1.5cm}|L{4.3cm}|Y|}
\hline
\filacab \cab{ID} & \cab{Decisión} & \cab{Qué impone a esta iteración} \\ \hline
\endhead
DEC-01 & El sistema sigue un modelo cliente-servidor basado en servicios. &
La descomposición del servidor produce servicios con responsabilidades separadas, no un solo proceso monolítico que lo resuelva todo. El cliente es un consumidor de esos servicios. \\ \hline
DEC-02 & No se confía en el cliente: toda la lógica del juego se ejecuta en el servidor. &
El cliente no valida reglas ni calcula resultados, ni siquiera para mostrarlos. Lo que la interfaz necesita saber para dibujar (casillas legales, surcos válidos, efecto de un muro sobre el camino más corto) se lo tiene que dar el servidor. \\ \hline
DEC-03 & Una partida en línea exige conexión con la base de datos. &
La partida no se juega sin base de datos disponible: PRE-04 de CU-20, CU-21, CU-25 y CU-26 se mantiene. Ante una caída, la partida no continúa a ciegas, sino que se suspende de forma controlada y se reanuda desde lo persistido. \\ \hline
\end{tabla}

\subsubsection*{Concerns}

\begin{tabla}{|L{1.5cm}|Y|}
\hline
\filacab \cab{ID} & \cab{Concern y qué aporta a la iteración} \\ \hline
\endhead
CRN-10 & \emph{Tengo que decidir quién manda sobre el estado de la partida, el cliente o el servidor. De esa decisión cuelga casi todo lo demás, y cambiarla más adelante me sale carísimo.} Es el concern que da origen al driver. \\ \hline
CRN-23 & El especialista da por hecho que alguien modificará el cliente y probará con uno alterado. Fija el supuesto de confianza con el que se diseña. \\ \hline
CRN-06 & El jugador quiere saber qué pasa si se le corta la conexión: si puede volver, si pierde o si deja al otro esperando. \\ \hline
CRN-19 & El administrador de base de datos exige recuperar cuentas, progreso y partidas, y que se defina qué información es aceptable perder. \\ \hline
CRN-02 & El jugador espera respuesta inmediata al colocar una barrera. Es el contrapeso de todo lo anterior. \\ \hline
CRN-11 & Sin WebSockets hay que definir dónde empieza y termina cada mensaje y qué hacer cuando cliente y servidor no tienen la misma versión. \\ \hline
\end{tabla}

\subsubsection*{Preguntas abiertas que condicionan la iteración}

\begin{tabla}{|L{1.3cm}|Y|}
\hline
\filacab \cab{ID} & \cab{Cómo condiciona} \\ \hline
\endhead
Q-01 & Si habrá un oponente controlado por IA. Afecta a quién ocupa el lugar del segundo jugador, no a quién tiene la autoridad. La estructura elegida no debe impedir que la IA entre como un participante más. \\ \hline
Q-02 & Si habrá conexión LAN. No cambia la autoridad, pero obliga a que el servidor de partidas sea un proceso desplegable por separado y a que su dirección se lea de configuración. \\ \hline
Q-09 & Con qué mecanismo se cifra el tráfico. Su primera mitad, que sí se cifra, ya está resuelta por ASR-03; el mecanismo se decide en la iteración~2 y cargará un costo al presupuesto de ESC-01. \\ \hline
\end{tabla}

% ------------------------------------------------------------
\subsection{Paso 2. Objetivo de la iteración y selección de entradas}

\subsubsection*{Priorización de las entradas}

Cada entrada del paso 1 se pesa por lo que aporta a este objetivo y no por su importancia general: una entrada puede ser crítica para el proyecto y pesar poco aquí si cualquier alternativa razonable la cumple.

\begin{tabla}{|L{3.2cm}|L{1.3cm}|Y|L{2.2cm}|}
\hline
\filacab \cab{Entrada} & \cab{Peso} & \cab{Por qué pesa lo que pesa} & \cab{Decisión} \\ \hline
\endhead
ESC-05 (ASR-02) & Alto & Es la entrada que decide la autoridad: con el cliente decidiendo, la medida es imposible; con el servidor decidiendo, es alcanzable. & Se atiende \\ \hline
ESC-08 (ASR-04) & Alto, pero de otro elemento & Obliga a que el estado sobreviva a la conexión que lo transportaba. Necesita decidir sobre la sesión de partida y sobre el protocolo, que todavía no existe. & Se difiere a una iteración posterior \\ \hline
ESC-09 (ASR-05) & Alto, pero de otro elemento & Obliga a definir qué ocurre cuando la base de datos se cae y qué se pierde. Necesita decidir sobre el servicio de acceso a datos. & Se difiere a la iteración~5 \\ \hline
CRN-10 & Alto & Enuncia la decisión y advierte su costo de cambio. Es el motivo por el que la iteración existe. & Se atiende \\ \hline
CU-20, CU-21 & Alto & Son las dos acciones que cambian el estado. Sus comprobaciones definen qué tiene que saber el árbitro en el momento de decidir. & Se atiende \\ \hline
CU-26 & Medio & Es el caso que exige separar sesión y conexión, que es el objetivo de una iteración posterior. Aquí solo se respeta la condición de no atar el estado al socket. & Se difiere a una iteración posterior \\ \hline
CU-25 & Medio & Define la frontera entre lo que se puede perder y lo que no, que es lo que mide ESC-09. Aquí solo aporta el estado final de la partida a la máquina de estados. & Se difiere a la iteración~5 \\ \hline
CON-02, CON-03, CON-17 & Alto & Determinan que la caída de una conexión no se note por sí sola y que la detección haya que construirla. & Se atiende \\ \hline
DEC-01, DEC-02 & Alto & Fijan el estilo y el reparto de responsabilidades: servicios en el servidor y un cliente sin lógica de juego. Con ellas, el elemento a descomponer deja de admitir alternativas repartidas entre los dos extremos. & Se atiende \\ \hline
DEC-03 & Alto & Convierte la disponibilidad de la base de datos en condición de la partida y decide qué pasa cuando se pierde: suspender en lugar de seguir arbitrando. & Se atiende \\ \hline
ESC-01 (ASR-01) & Alto, como límite & No es el objetivo, pero descarta toda alternativa que agregue viajes por jugada o que espere a un tercero antes de responder. Con DEC-02, la respuesta cuesta un viaje completo de red, así que el margen del servidor se estrecha. & Se atiende como restricción \\ \hline
CU-24 & Bajo & El plazo de los 60 segundos para reclamar la victoria depende de la desconexión, que se trata en una iteración posterior. & Se difiere a una iteración posterior \\ \hline
CRN-23 & Medio & Es el origen de ESC-05: fija el supuesto de que el cliente es hostil, que es lo que obliga a poner la autoridad en el servidor. & Se atiende \\ \hline
CRN-06, CRN-19 & Medio & Son el origen de ESC-08 y ESC-09; se atienden a través de sus medidas, en las iteraciones que las toman. & Se difiere \\ \hline
ESC-07, CRN-02 & Medio & ESC-07 pide un registro que ESC-05 ya exige; CRN-02 llega como la medida de ESC-01. & Se atiende por vía indirecta \\ \hline
CON-04, DEP-03, DEP-04 & Medio & Fijan qué puede fallar, no cómo se estructura la respuesta a la falla. & Se atiende como supuesto \\ \hline
CU-17, CU-18, CU-19 & Bajo & Crean la partida, pero cualquier forma de crearla es compatible con las alternativas en juego. Aquí solo importa que, cuando la partida nace, su estado ya tiene dueño. & Se difiere al servicio de emparejamiento y salas \\ \hline
CRN-11, ASR-03 (ESC-06) & Bajo aquí & Delimitación, versionado y cifrado son del protocolo: condicionan el costo por mensaje, no quién manda. & Se difiere a la iteración~2 \\ \hline
ASR-06 (ESC-12) & Bajo aquí & Con DEC-02 el motor tiene un solo consumidor, el servidor. Esta iteración solo necesita que exista un motor completo al que el arbitraje consulte; que sea parametrizable y ejecutable sin interfaz es otro objetivo. & Se difiere a la iteración~3 \\ \hline
Q-01, Q-02 & Bajo & Ninguna cambia la autoridad. Solo imponen que la IA pueda entrar como un participante más y que el servidor sea desplegable por separado. & Se atiende como condición \\ \hline
CON-13, CON-14 & Bajo & No eligen entre alternativas; descartan las que no se puedan construir por partes. & Se atiende como filtro \\ \hline
\end{tabla}

\subsubsection*{Objetivo de la iteración}

\begin{tabla}{|L{3.0cm}|Y|}
\fichacab
\parte{Objetivo}{Que el servidor sea la única autoridad sobre el estado de una partida en curso: ningún mensaje de un cliente, alterado o no, cambia el estado si no es una jugada legal del participante a quien le toca el turno.}
\parte{ASR que satisface}{ASR-02, sobre la base de ASR-08: la partida en línea arbitrada es lo que el incremento recorre, y la autoridad del servidor es lo que la medida comprueba.}
\parte{Driver que cubre}{DRV-01, en lo que hace a la autoridad y al estado vigente.}
\parte{Medida}{La de ESC-05: con 200 mensajes inválidos de un cliente modificado, 50 al azar, el 100~\% se rechaza, el estado es idéntico antes y después, 0 respuestas llevan datos del rival y el 100~\% de los intentos queda registrado con cuenta y hora.}
\parte{Límite que respeta}{ESC-01. Ninguna decisión de esta iteración puede impedir que una jugada se responda en menos de un segundo; el reparto del presupuesto se decide en la iteración~4.}
\parte{Incremento que produce}{Dos clientes juegan una partida completa en línea contra un servidor que arbitra: CU-20 y CU-21 funcionan de principio a fin, el cliente resalta lo que el servidor le indica y un cliente de prueba que envía mensajes inválidos no consigue alterar la partida. Es demostrable al cerrar la iteración y sirve como incremento semanal (CON-14).}
\parte{Criterio de éxito}{La iteración termina cuando ese incremento pasa la medida de ESC-05 sin violar ESC-01. Una alternativa de diseño que no lo permita se descarta, aunque sea mejor en lo demás.}
\end{tabla}


% ------------------------------------------------------------
\subsection{Paso 3. Elemento a descomponer}

Este es el elemento sobre el que trabaja la iteración. La ficha dice qué abarca, por qué se eligió frente a otros recortes posibles del sistema, qué queda por decidir sobre él y contra qué elementos del análisis se puede rastrear.

% ------------------------------------------------------------
\Needspace*{0.4\textheight}
\subsubsection*{EL-01 --- Gestión del estado de la partida}

\begin{tabla}{|L{2.6cm}|Y|}
\fichacab
\parte{Elemento}{La gestión del estado de la partida, entendida como un servicio del servidor (DEC-01) que concentra toda la lógica del juego (DEC-02).}
\parte{Qué abarca}{El estado vigente de cada partida en curso: el tablero con sus muros, la posición de cada peón, los muros restantes de cada participación, el turno, el número de jugada, los relojes y el estado de la partida. Y las operaciones que lo cambian: las dos acciones del turno (CU-20 y CU-21) y las condiciones de término (CU-22 a CU-25).}
\parte{Por qué se elige}{No por ser un módulo del sistema, sino porque la medida del objetivo depende de él: ESC-05 exige que el estado de la partida sea idéntico antes y después de 200 mensajes inválidos, y eso solo se puede afirmar de algo que tenga un dueño único y conocido. Se elige frente a descomponer el sistema completo o el servidor entero, que arrastrarían cuentas, tienda, chat y emparejamiento a una decisión de la que ninguno depende; y frente a descomponer el cliente, que es consecuencia y no causa, porque mientras no se sepa quién decide una jugada no se puede decir qué le toca a él.}
\parte{Qué falta decidir sobre él}{Cómo se mantiene: no está establecido si el estado vigente vive en la memoria del servidor, en la base de datos o repartido entre los clientes. Cómo se produce cada cambio: no está establecido quién comprueba el turno, el número de jugada y la legalidad antes de que el estado cambie, ni qué ocurre si dos mensajes llegan a la vez para la misma partida. Esas decisiones cambian a la vez el cliente, el servidor y el protocolo, porque determinan si un mensaje es una propuesta o un hecho consumado, y son las más caras de revertir (CRN-10).}
\parte{Evidencia en los casos de uso}{CU-21 reconstruye el tablero a partir de las Jugada registradas (paso 12) y confirma una transacción antes de notificar (pasos 15 a 19), sin decir de dónde sale el tablero con el que valida. CU-26 recupera todas las Jugada desde la base de datos al reconectar (paso 9), lo que solo es necesario si el estado no está ya en memoria. CU-20, CU-21, CU-25 y CU-26 exigen en PRE-04 conexión disponible con la base de datos, que DEC-03 mantiene. Y CU-21 resalta los surcos válidos (paso 3) y muestra el efecto del muro sobre el camino más corto (paso 5), cálculos que con DEC-02 tiene que entregar este servicio.}
\parte{Trazabilidad}{DRV-01. DEC-01, DEC-02, DEC-03. ESC-05, con ESC-01 como límite. CU-20, CU-21, CU-22, CU-23, CU-24, CU-25, CU-26. REQ-01, CON-04. CRN-10, CRN-23, CRN-02. STK-05, STK-15, STK-01.}
\end{tabla}

% ------------------------------------------------------------
\subsection{Paso 4. Conceptos de diseño}

Estos son los conceptos de diseño con los que se descompone la gestión del estado de la partida. Cada ficha dice qué atiende y por qué se eligió frente a las alternativas, qué efectos trae consigo y contra qué elementos del análisis se puede rastrear.

% ------------------------------------------------------------
\Needspace*{0.3\textheight}
\subsubsection*{CD-01 --- Cliente-servidor basado en servicios, con servidor autoritativo}

\begin{tabla}{|L{2.6cm}|Y|}
\fichacab
\parte{Estilo}{Estilo arquitectónico.}
\parte{Descripción}{El cliente pide y presenta; el servidor decide. Atiende ESC-05 porque un mensaje del cliente es una propuesta y no un hecho: lo que cambia el estado es la decisión del servidor. Se elige frente al cliente autoritativo con servidor retransmisor y frente a la arquitectura entre pares, que cuestan menos latencia pero dejan el estado en manos de un cliente que se supone hostil (CRN-23). El estilo, además, ya está fijado por DEC-01.}
\parte{Efectos}{Cada jugada cuesta un viaje completo de red, que consume presupuesto de ESC-01. El servidor necesita el motor de reglas completo. El cliente queda reducido a presentación y entrada, y su versión deja de ser relevante para la corrección de la partida.}
\parte{Trazabilidad}{DRV-01. DEC-01, DEC-02. ESC-05 (ASR-02), con ESC-01 como límite. CU-20, CU-21. REQ-01. CRN-10, CRN-23. STK-05, STK-15, STK-10.}
\end{tabla}

% ------------------------------------------------------------
\Needspace*{0.3\textheight}
\subsubsection*{CD-02 --- Toda la lógica del juego en el servidor}

\begin{tabla}{|L{2.6cm}|Y|}
\fichacab
\parte{Estilo}{Estilo: reparto de responsabilidades entre los dos extremos.}
\parte{Descripción}{El cliente no calcula reglas ni resultados, ni siquiera para mostrarlos. Atiende ESC-05 en su raíz: si el cliente no calcula nada, no hay cálculo suyo que el servidor tenga que desconfiar o contrastar. Se elige frente a duplicar el motor en el cliente para previsualizar, que es lo habitual en juegos en línea y lo que DEC-02 descarta.}
\parte{Efectos}{La notificación de turno tiene que llevar lo que la interfaz necesita dibujar, que es de donde sale CD-11. Descarta la validación optimista y estrecha el margen de ESC-01, porque ya no hay respuesta local que adelantar. A cambio, el motor de reglas tendrá un solo consumidor, lo que simplifica la iteración~3.}
\parte{Trazabilidad}{DEC-02. DRV-01, DRV-05. ESC-05, ESC-01. CU-20 (pasos 1 a 3), CU-21 (pasos 3 a 7). CRN-23, CRN-12. STK-15, STK-05.}
\end{tabla}

% ------------------------------------------------------------
\Needspace*{0.3\textheight}
\subsubsection*{CD-03 --- Servicios separados dentro del servidor}

\begin{tabla}{|L{2.6cm}|Y|}
\fichacab
\parte{Estilo}{Estilo: descomposición en servicios.}
\parte{Descripción}{El servidor se divide en protocolo y sesiones, arbitraje de partida, y acceso a datos. Atiende el objetivo al darle al estado de la partida un dueño con frontera propia, separado de lo que lo transporta y de lo que lo guarda. Se elige frente a un servidor de una sola pieza, que mezclaría las tres cosas y obligaría a decidirlas a la vez, en contra de CON-14.}
\parte{Efectos}{Fija las fronteras de las iteraciones siguientes: protocolo en la~2, sesión en la~3, datos en la~6. Obliga a definir interfaces internas antes de construir. Cada salto entre servicios cuesta tiempo dentro del segundo de ESC-01, así que su número se mantiene bajo.}
\parte{Trazabilidad}{DEC-01. DRV-01, DRV-02, DRV-04, DRV-08. ESC-05, ESC-01. CON-03, CON-13, CON-14. CRN-11. STK-05, STK-14.}
\end{tabla}

% ------------------------------------------------------------
\Needspace*{0.3\textheight}
\subsubsection*{CD-04 --- Una máquina de estados por partida}

\begin{tabla}{|L{2.6cm}|Y|}
\fichacab
\parte{Estilo}{Patrón de diseño.}
\parte{Descripción}{La partida tiene estados y transiciones explícitos, con EN\_CURSO y FINALIZADA junto a su forma de término (RN-01 del CU-25). Atiende ESC-05 porque un mensaje que no corresponde al estado vigente se rechaza por definición y no por una comprobación suelta, y porque toda condición de término pasa por el mismo lugar. Se elige frente a resolver cada final en el caso de uso que lo provoca, que abriría cinco caminos distintos para cerrar la misma partida.}
\parte{Efectos}{El reloj del turno y los plazos viven en el servidor, dentro de la máquina de estados, no en el cliente. CU-22, CU-23, CU-24 y CU-25 convergen en un único punto de cierre. Obliga a enumerar las transiciones antes de codificar, que es trabajo de diseño por adelantado.}
\parte{Trazabilidad}{DRV-01. ESC-05. CU-20, CU-21, CU-22, CU-23, CU-24, CU-25 (RN-01). CRN-10. STK-05, STK-01.}
\end{tabla}

% ------------------------------------------------------------
\Needspace*{0.3\textheight}
\subsubsection*{CD-05 --- Objeto activo por partida, con cola de mensajes}

\begin{tabla}{|L{2.6cm}|Y|}
\fichacab
\parte{Estilo}{Patrón de concurrencia.}
\parte{Descripción}{Cada partida atiende sus mensajes de uno en uno, en el orden en que llegan a su cola. Atiende ESC-05 en la parte que exige que el estado sea idéntico antes y después: sin serialización, dos mensajes simultáneos podrían dejar la partida en un estado que ninguna jugada legal produce. Además hace que una jugada fuera de orden se detecte de forma natural (FA-07 del CU-21). Se elige frente a un bloqueo global compartido, que serializaría entre sí las 20 partidas y rompería ESC-01.}
\parte{Efectos}{Cada partida avanza con independencia de las demás, lo que permite medir sus tiempos por separado. El trabajo por jugada tiene que ser corto, porque mientras dura bloquea a su propia cola. El estado de una partida no se toca desde fuera de su objeto activo.}
\parte{Trazabilidad}{DRV-01, DRV-03. ESC-05, ESC-01. CU-20, CU-21 (FA-07, NOT\_YOUR\_TURN). CRN-02, CRN-10. STK-01, STK-15.}
\end{tabla}

% ------------------------------------------------------------
\Needspace*{0.3\textheight}
\subsubsection*{CD-06 --- Copia de trabajo en memoria con respaldo obligatorio en la base de datos}

\begin{tabla}{|L{2.6cm}|Y|}
\fichacab
\parte{Estilo}{Patrón de gestión de estado.}
\parte{Descripción}{El servicio valida y responde contra la partida que tiene en memoria, y cada jugada se escribe en SQL Server, que es la copia de referencia de la que se reconstruye. Responde a la pregunta del objetivo, dónde vive el estado vigente, y es lo que permite afirmar que el estado sea idéntico antes y después de 200 mensajes inválidos (ESC-05): sin un dueño único no hay nada de lo que afirmarlo. Se elige frente al servicio sin estado que lee la partida completa en cada jugada, que metería la latencia de DEP-03 y DEP-04 en cada turno.}
\parte{Efectos}{Cumple DEC-03 sin poner la base de datos en el camino de la respuesta. Obliga a reconstruir desde las jugadas registradas al recuperar, que es lo que CU-21 (paso 12) y CU-26 (paso 9) ya describen. Deja abierta, para la iteración~5, la pregunta de qué ocurre cuando esa base de datos no está.}
\parte{Trazabilidad}{DEC-03. DRV-01, DRV-04. ESC-05, ESC-01. CU-20, CU-21 (paso 12), CU-26 (paso 9). CON-04, DEP-03, DEP-04. CRN-19, CRN-10. STK-13, STK-14.}
\end{tabla}

% ------------------------------------------------------------
\Needspace*{0.3\textheight}
\subsubsection*{CD-07 --- Validación de la entrada antes de que llegue al motor}

\begin{tabla}{|L{2.6cm}|Y|}
\fichacab
\parte{Estilo}{Táctica de seguridad: resistir ataques, validar la entrada.}
\parte{Descripción}{La delimitación, la estructura y la versión del mensaje se comprueban antes de que este toque el estado. Atiende ESC-05 en la parte de los mensajes mal formados, truncados o de otra versión del protocolo, que son 50 de los 200 de la medida. Se elige frente a dejar que el arbitraje se defienda por su cuenta, que repartiría la misma comprobación por todo el servicio.}
\parte{Efectos}{El arbitraje solo recibe mensajes bien formados y puede ocuparse de reglas y turno. El mecanismo concreto pertenece al protocolo y se decide en la iteración~2. Agrega un paso por mensaje, con su costo dentro de ESC-01.}
\parte{Trazabilidad}{DRV-01, DRV-02. ESC-05. CU-21 (EX-02, EX-03). CON-03. CRN-11, CRN-23. STK-15, STK-10, STK-05.}
\end{tabla}

% ------------------------------------------------------------
\Needspace*{0.3\textheight}
\subsubsection*{CD-08 --- Autenticación y autorización de cada mensaje}

\begin{tabla}{|L{2.6cm}|Y|}
\fichacab
\parte{Estilo}{Táctica de seguridad: autenticar y autorizar actores.}
\parte{Descripción}{Cada mensaje se comprueba contra una sesión abierta y contra la participación de esa cuenta en esa partida, tal como CU-21 hace en su paso 8 y resuelve en FA-05 con SESSION\_INVALID y NOT\_A\_PARTICIPANT. Atiende ESC-05 en la parte que exige 0 respuestas con datos del rival: sin esta comprobación, un cliente alterado podría pedir el estado de una partida ajena. Se elige frente a confiar en el identificador de jugador que venga dentro del mensaje, que es precisamente lo que un cliente alterado falsifica.}
\parte{Efectos}{El estado de la partida necesita conocer la sesión de cada participante, pero no su conexión, lo que deja preparada la separación de una iteración posterior. Obliga a que el módulo de cuentas exponga la validez de una sesión al servicio de partida.}
\parte{Trazabilidad}{DRV-01, DRV-06. ESC-05. CU-21 (paso 8, FA-05), CU-20, CU-01. CRN-23, CRN-05. STK-15, STK-02.}
\end{tabla}

% ------------------------------------------------------------
\Needspace*{0.3\textheight}
\subsubsection*{CD-09 --- Registro de cada intento inválido con cuenta y hora}

\begin{tabla}{|L{2.6cm}|Y|}
\fichacab
\parte{Estilo}{Táctica de seguridad: detectar y auditar.}
\parte{Descripción}{Todo mensaje rechazado queda registrado con la cuenta que lo envió y el momento en que llegó. Atiende la última parte de la medida de ESC-05, que exige el 100~\% de los intentos registrados, y es lo que convierte un rechazo silencioso en evidencia. Se elige frente a contar los rechazos sin identificarlos, que no permitiría investigar ni sostener una sanción.}
\parte{Efectos}{Exige que todo mensaje llegue ligado a una cuenta, lo que depende de CD-08. Genera escritura adicional, que por CD-10 no puede entrar en el camino de la respuesta. Da la base sobre la que después se apoyan la moderación de CU-42 a CU-45 y las bitácoras de CU-48.}
\parte{Trazabilidad}{DRV-01, DRV-06. ESC-05, ESC-07. QA-08. CU-42, CU-43, CU-44, CU-48. CRN-18, CRN-23. STK-15, STK-16, STK-17.}
\end{tabla}

% ------------------------------------------------------------
\Needspace*{0.3\textheight}
\subsubsection*{CD-10 --- Escritura fuera del camino de la respuesta}

\begin{tabla}{|L{2.6cm}|Y|}
\fichacab
\parte{Estilo}{Táctica de desempeño: limitar el trabajo dentro de la respuesta.}
\parte{Descripción}{El servicio valida, responde y después escribe. Atiende el límite de ESC-01 sin contradecir DEC-03: lo que esa decisión exige es que la conexión con la base de datos exista, no que el jugador espere a que la transacción confirme. Se elige frente a confirmar la transacción antes de notificar, como CU-21 describe hoy en sus pasos 15 a 19, que suma a cada turno la latencia de una infraestructura externa.}
\parte{Efectos}{Abre una ventana entre lo respondido y lo escrito, que ESC-09 acota en cinco segundos y que la iteración~5 define. Obliga a una cola de escrituras con constancia de lo pendiente. Si la iteración~4 muestra que la confirmación cabe en el presupuesto, la ventana se acorta sin cambiar la estructura.}
\parte{Trazabilidad}{DEC-03. DRV-01, DRV-03, DRV-04. ESC-01, ESC-09. CU-20, CU-21 (pasos 15 a 19), CU-25. TEN-06. CRN-02, CRN-19. STK-01, STK-13.}
\end{tabla}

% ------------------------------------------------------------
\Needspace*{0.3\textheight}
\subsubsection*{CD-11 --- El servidor entrega las jugadas legales junto con el turno}

\begin{tabla}{|L{2.6cm}|Y|}
\fichacab
\parte{Estilo}{Táctica de desempeño: limitar el trabajo, con origen en DEC-02.}
\parte{Descripción}{Con la notificación de turno viajan las casillas legales, los surcos válidos y el efecto de cada muro sobre el camino más corto. Es lo que sustituye a la validación previa en el cliente que DEC-02 descarta: la interfaz resalta y previsualiza con datos que ya recibió, sin calcular nada y sin volver a consultar mientras el jugador decide. Se elige frente a preguntar al servidor cada vez que el jugador mueve el muro por el tablero, que multiplicaría los viajes dentro del turno y rompería ESC-01.}
\parte{Efectos}{Aumenta el tamaño del mensaje de turno, con su costo para el protocolo de la iteración~2 y para el presupuesto de la~4. Mantiene CU-20 y CU-21 tal como están escritos de cara al jugador. Obliga a que el motor de reglas sepa enumerar las jugadas legales y explicar por qué una no lo es, que es una condición para la iteración~3 y para ESC-02.}
\parte{Trazabilidad}{DEC-02. DRV-01, DRV-05. ESC-01, ESC-02, ESC-05. CU-20 (pasos 1 a 3), CU-21 (pasos 3 a 5, FA-02). REQ-07. CRN-01, CRN-02. STK-01, STK-09.}
\end{tabla}
